% =====================================================================
% project_report.tex - Work-submission report
% Evaluating India's 2021 Agricultural Derivatives Suspension
% Compile with pdflatex (no special dependencies beyond standard packages).
% =====================================================================
\documentclass[11pt]{article}
\usepackage[a4paper,margin=1in]{geometry}
\usepackage{amsmath,amssymb}
\usepackage{booktabs}
\usepackage{enumitem}
\usepackage{framed}
\usepackage[hidelinks]{hyperref}
\hypersetup{pdfauthor={Naman Karwa}}
\setlength{\parskip}{4pt}

\newenvironment{keybox}
  {\begin{framed}\noindent\textbf{Key point.}\ }
  {\end{framed}}

\title{\textbf{Did India's 2021 Suspension of Agricultural Derivatives\\
Deliver Price Stability?}\\[6pt]
\large A Complete Work Report: Data, Methods, Errors, and Findings}
\author{Naman Karwa\\ Student, Department of Mathematics\\ Indian Institute of Technology, Guwahati}
\date{July 2026}

\begin{document}
\maketitle

\begin{abstract}
\noindent In December 2021 India's securities regulator (SEBI) suspended exchange-traded
derivatives in seven agricultural commodities, on the premise that futures speculation was
driving up food prices. This report documents, end to end, an empirical evaluation of that
policy: how every dataset was collected and cleaned, how the research design evolved, the
errors I made and how each was caught and corrected, and where the findings currently
stand. The short version of the result: the fear that switching off futures would make
cash (spot) markets \emph{more} volatile is not borne out - if anything spot volatility
was lower afterwards - but the size of that calming effect depends heavily on which
comparison commodities are used, and under the most economically appropriate comparison
group the aggregate effect is modest (about $-10\%$) and not statistically significant,
with the only strong individual case being chana (chickpea). Companion analyses suggest
the real cost of the suspension lies in price discovery and market integration rather than
in lost hedging. I present this as an honest, quasi-experimental evaluation - suggestive,
not conclusive - together with everything needed to check, reproduce, and extend it.
\end{abstract}

\tableofcontents
\newpage

% =====================================================================
\section{Introduction}

Food inflation is politically sensitive in India, and futures markets are a recurring
suspect: the claim is that speculation in derivatives inflates the underlying spot
(``mandi'', i.e.\ wholesale market) price. Economic theory is genuinely ambivalent -
futures can \emph{stabilise} prices by aggregating information and enabling hedging, or
can be blamed for destabilising them - and India's policy history reflects that
ambivalence: commodity futures have been banned and re-permitted repeatedly since the
1950s, and the 2008 Abhijit Sen Committee found no conclusive evidence that futures raise
spot volatility.

On 19-20 December 2021, SEBI suspended derivatives trading in seven agricultural
commodities - wheat, paddy (non-basmati), chana (chickpea), mustard seed and its
complex, soybean and its complex, crude palm oil, and moong - initially for one year,
and subsequently extended the suspension repeatedly (at the time of writing, to March
2027). Two commodities were hit earlier: new positions in chana were barred from
16~August~2021, and in mustard from 8~October~2021. This is a sharp, dated policy
intervention that switched off futures for a specific set of commodities while leaving
others trading - a natural experiment for the question:

\begin{keybox}
Did the suspension deliver what it promised - calmer spot prices - and what did it
cost in market function? Concretely: how did the \emph{spot-price volatility} of the
banned commodities change after the suspension, relative to comparable commodities whose
futures kept trading?
\end{keybox}

This report is a complete account of the work: Section~\ref{sec:design} sets out the
research design; Section~\ref{sec:data} documents the data collection in detail (this was
the hardest part of the project and the part I document most carefully);
Section~\ref{sec:cleaning} covers cleaning and quality control;
Section~\ref{sec:errors} is an explicit register of the errors I made and how each was
caught and fixed; Section~\ref{sec:methods} explains the methods in plain language;
Section~\ref{sec:progression} traces how the findings evolved as the data and methods
improved - the evolution is itself a finding; Section~\ref{sec:results} presents the
current results; Section~\ref{sec:limits} states the limitations honestly; and
Section~\ref{sec:repro} explains how to navigate the folder and reproduce every number.

% =====================================================================
\section{Research design at a glance}
\label{sec:design}

\paragraph{Treated and comparison commodities.} The seven suspended commodities form the
treatment group. Commodities whose futures \emph{kept trading} form the comparison
(``donor'') group: initially castor seed, guar seed, jeera (cumin), turmeric, and cotton;
later extended - for reasons explained in Section~\ref{sec:progression} - with
non-banned \emph{food} crops (barley, maize, jowar, bajra, with groundnut, sesamum,
sunflower and ragi in progress).

\paragraph{Outcome.} Spot-price volatility: the 30-day rolling standard deviation of
daily log returns of mandi prices, annualised (``realized volatility''), computed at the
commodity-district level and aggregated monthly. Spot prices are the right outcome
because the banned commodities have no futures prices after the suspension, and because
``spot stability'' is precisely what the policy promised.

\paragraph{Design principles fixed in advance.}
\begin{itemize}[leftmargin=1.6em,itemsep=1pt]
  \item \emph{Raw data is immutable}: every downloaded file is stored verbatim and never
        edited; all cleaning happens in scripts, so every transformation is inspectable.
  \item \emph{Every acquisition is logged}: one row per data series in a provenance log
        (source, dates, units, known issues).
  \item \emph{Pre-registered verification}: the volatility measure, the estimating
        equations, and explicit pass/fail rules (including falsification tests the design
        must survive) were written down \emph{before} estimation.
  \item \emph{Banned-commodity futures are hard-truncated} at 20~December~2021 in all
        cleaned files: exchanges kept publishing stale ``settlement'' quotes for months
        after the suspension with zero trading, and treating those as prices would
        fabricate post-ban futures data.
  \item \emph{Methodological decisions are dated and documented} with their rationale
        (the full decision history, including choices later revised, is laid out in
        Sections~\ref{sec:errors} and~\ref{sec:progression} of this report), so the path
        from first idea to final specification is auditable.
\end{itemize}

% =====================================================================
\section{Data: what was collected and how}
\label{sec:data}

Getting the data was the single hardest part of this project. No single source provides
what the design needs, several official archives turned out to be unusable for the
required window, and every workaround forced a documented decision. Table~\ref{tab:data}
summarises the holdings; the subsections then describe, source by source, how each was
obtained.

\begin{table}[htbp]
\centering\small
\caption{Data holdings at a glance.}
\label{tab:data}
\begin{tabular}{p{3.6cm} p{3.4cm} p{2.2cm} p{5.2cm}}
\toprule
Series & Source & Period & Volume / notes \\
\midrule
Mandi (spot) prices, national & CEDA data portal (Ashoka University), re-serving Govt.\ of India Agmarknet & 2017-Oct 2025 & 20 commodities $\times$ $\sim$3{,}200 daily rows each \\
Mandi (spot) prices, district & same & 2017-Oct 2025 & $\sim$5.74 million market-day rows (12 commodities) $+$ $\sim$1.2 million (food-donor extension, ongoing) \\
Futures, comparison commodities & Investing.com (vendor) & 2017-2026 & castor, guar, jeera, turmeric, kapas; three contract legs each (c1/c2/c3) \\
Futures, crude palm oil (banned) & MCX exchange archive & 2017-2021 & all 64 monthly contracts, with volume and open interest \\
Futures, cotton & MCX exchange archive & 2017-2026 & both contract generations (77 $+$ 15 contracts) \\
Participant composition & SEBI monthly bulletins & 2019-2021 & 67 files (33 of 36 months); participant-wise turnover shares \\
International references & public financial-data API & 2017-2025 & CBOT wheat; USD/MYR exchange rate \\
Policy timeline & primary documents (SEBI, DGFT, PIB, CBIC) & 2021-2026 & 24 dated interventions, each with a source URL \\
\bottomrule
\end{tabular}
\end{table}

\subsection{Spot (mandi) prices - the core dataset}
\label{sec:data-spot}

\paragraph{Why it was hard.} India's official wholesale-price portal (Agmarknet) exposes
prices mandi-by-mandi through web forms that are impractical for a nine-year,
multi-hundred-district panel. I found that the Centre for Economic Data and Analysis
(CEDA) at Ashoka University re-serves the same Ministry of Agriculture Agmarknet data,
cleaned, through a documented programmatic interface (free academic registration). I chose
that route and cite both CEDA and the underlying Agmarknet source; the data are
Rs-per-quintal modal (most frequent) transaction prices per market per day.

\paragraph{Traps and how they were handled.}
\begin{itemize}[leftmargin=1.6em,itemsep=1pt]
  \item \emph{Rate limit.} The interface allows only 40 requests per hour. A naive
        first attempt silently exhausted the quota (17 minutes, zero rows). I rebuilt the
        downloader to read the server's rate-limit headers and sleep to the reset window,
        and made it fully resumable (one file per unit of work, finished files skipped),
        so a multi-hour pull survives interruptions. The full district panel -
        396 state-files, 5.74~million market-day rows - took roughly ten hours of paced
        requests.
  \item \emph{Oversized responses.} Large states (e.g.\ Uttar Pradesh) time out when
        requested whole. The downloader batches districts and, when a single dense
        district still fails, recursively halves the date range until the server returns
        it.
  \item \emph{Commodity-identifier mismatch.} The CEDA commodity identifiers do not match
        Agmarknet's own identifiers (e.g.\ turmeric is 39 in one and 35 in the other). I
        verified every commodity identifier against the source's commodity list before
        pulling.
  \item \emph{A mislabelled series.} Cross-checking every spot series against its futures
        counterpart exposed that the series labelled ``Guar'' (id~75) was contaminated
        with guar \emph{gum} - a processed product roughly ten times the price of guar
        \emph{seed}. Its correlation with guar-seed futures was 0.04. The separate series
        ``Guar Seed'' (id~413) correlates 0.99 with the futures and replaced it. This
        catch is a good illustration of why every series was validated against an
        independent benchmark before use.
\end{itemize}

\paragraph{The food-donor extension.} After the first round of results (see
Section~\ref{sec:progression}), I extended the spot pull to eight non-banned \emph{food}
commodities - barley, maize, jowar, bajra, ragi, groundnut, sesamum, sunflower - to
serve as economically comparable controls. At the time of writing the cereals (barley,
maize, jowar, bajra) are fully or nearly fully downloaded and are \emph{in} the analysis
panel; the oilseeds and ragi are still downloading under the 40-per-hour limit.

\subsection{Futures prices}
\label{sec:data-futures}

\paragraph{The archive wall.} The natural source, the NCDEX exchange's public daily-price
archive, only reaches back to mid-2024 - useless for a 2017-2021 pre-period. Three
workarounds were used, and one gap remains:

\begin{itemize}[leftmargin=1.6em,itemsep=1pt]
  \item \emph{Comparison commodities via a vendor.} Investing.com serves daily
        open-high-low-close back to 2017 for the still-trading commodities. I downloaded
        each as three separate contract ``legs'' - front month (c1), next (c2), far
        (c3) - because liquidity and price behaviour differ along the curve. Because a
        continuous futures series is stitched from expiring contracts, I
        \emph{reverse-engineered the vendor's stitching rule} from the data itself
        (detecting days where the front series adopts the next contract's previous price):
        the vendor holds the front contract through its final trading day, switches the
        next session, and leaves the splice unadjusted. Knowing the exact convention let
        me replicate it when constructing series myself, so cross-commodity comparisons
        are not contaminated by inconsistent construction.
  \item \emph{Crude palm oil from the exchange itself.} Palm oil trades on MCX, whose
        website rejects all non-browser access. Working through a real browser session, I
        used the exchange's own data endpoint to pull \emph{all 64 monthly contracts}
        from 2017 to the suspension - including traded volume and open interest, which
        the vendor route lacks - and constructed continuous series replicating the
        verified vendor convention. A critical discovery in this raw data: contracts kept
        publishing settlement marks for months \emph{after} the suspension with zero
        volume. Those stale quotes are exactly why the cleaned data enforce a hard
        truncation at the ban date.
  \item \emph{Cotton, both generations.} MCX cotton changed contract specification
        (a bale-denominated contract until 2022, a candy-denominated one from 2023);
        both generations were pulled (77 $+$ 15 contracts). The unit break, and cotton's
        own trading halt in August 2022, are reasons cotton is treated cautiously as a
        control.
  \item \emph{The gap.} Pre-ban futures for the six NCDEX-listed banned commodities
        (wheat, chana, mustard, soybean, paddy, moong) are not publicly available for
        2017-2021. An academic data request to the exchange was submitted in June 2026
        and remains pending. This blocks the futures-based hypotheses
        (Section~\ref{sec:limits}); the spot-based analysis - the core of this report
 - does not depend on it.
\end{itemize}

\subsection{Regulator bulletins (who used these markets)}
SEBI's monthly bulletins carry participant-wise data for the commodity-derivatives
segment. The bulletin listing pages are dynamically rendered and unlinkable, but the
underlying document pages are served to standard web crawlers; I recovered the document
URLs through search and web-archive indexes and downloaded 33 of 36 months for
2019-2021 (67 files; three months are missing on the regulator's own server). The
spreadsheet tables in these bulletins are the input for the market-composition analysis
(Section~\ref{sec:results-companions}). A practical wrinkle documented in the code: the
relevant table's number and layout drift from month to month, so the parser locates it by
content, and 10 of 33 bulletins (the twelve months to March 2021) parse cleanly - enough
for a stable pre-ban composition estimate, with the coverage limitation stated wherever
the numbers are used.

\subsection{International references and the policy timeline}
CBOT wheat futures and the dollar-ringgit exchange rate (for palm-oil work) were pulled
from a public financial-data interface. Separately, I compiled a
\emph{policy-intervention timeline}: 24 dated events - the suspension itself, its six
extensions, export bans on wheat and rice, stock limits, edible-oil duty changes - each
verified against a primary document (SEBI press releases, DGFT notifications, gazette
orders) with the source URL recorded. This ledger matters because the suspension did not
happen in isolation: several banned commodities were simultaneously hit by other
interventions, which is a central confounding concern throughout.

% =====================================================================
\section{Data organisation, cleaning, and quality control}
\label{sec:cleaning}

\paragraph{Three tiers.} The data folder has three tiers with strict roles:
\texttt{raw/} (verbatim downloads, immutable), \texttt{constructed/} (deterministic
derivations from raw, e.g.\ continuous futures series, built by scripts),
and \texttt{clean/} (analysis-ready CSVs produced only by the cleaning scripts).
Nothing is ever hand-edited.

\paragraph{Cleaning rules.} The main ones, each enforced in code:
\begin{itemize}[leftmargin=1.6em,itemsep=1pt]
  \item \emph{Trading days only.} Mandi price feeds behave like seven-day calendar grids
 - across hundreds of districts, some market reports on virtually every calendar
        day, including weekends, often just carrying Friday's price. Returns are therefore
        computed on weekday observations only, so that the volatility annualisation
        (which assumes $\sqrt{252}$ trading days) matches the sampling grid. This rule
        exists because of an error I initially made; see Section~\ref{sec:errors}.
  \item \emph{Hard truncation of banned futures} at 20~December~2021 (stale post-ban
        settlement quotes are not prices).
  \item \emph{Exclusion of price-support-censored series.} For paddy, 40.3\% of daily
        returns are \emph{exactly zero} because government procurement at the minimum
        support price (MSP) pins the market price for long stretches; measured
        ``volatility'' of an administered price is a mechanical artifact. Paddy is
        excluded from the primary analysis; wheat (19.7\% zero returns) is retained but
        flagged, and results are shown with and without it.
  \item \emph{Independent benchmark checks.} Every spot series was correlated against its
        futures counterpart where one exists (this is what caught the guar mislabelling),
        and price levels, spans, gap structure, and extreme returns are printed as
        diagnostics every time a cleaning script runs.
\end{itemize}

\paragraph{Audit results.} A full audit of the cleaned data confirms: zero weekend
observations after the trading-day rule; no duplicate dates; no non-positive prices;
uniform 2017-October 2025 coverage across series ($\sim$2{,}200-2{,}300 trading days
each); and sane price levels. The main analysis panel is a commodity $\times$ district
$\times$ month panel of realized volatility: 172{,}381 cells across 16 commodities;
the headline regression below uses 146{,}507 observations across 2{,}244
commodity-district units (14 commodity clusters: 5 treated, 9 comparison).

% =====================================================================
\section{Errors made, and how each was caught and corrected}
\label{sec:errors}

Empirical work goes wrong in quiet ways, and I think the most useful thing I can document
is not that errors occurred but \emph{how the project's own checks caught them}.
Table~\ref{tab:errors} is the register; the two most consequential are narrated below.

\begin{table}[htbp]
\centering\small
\caption{Error register: every substantive error, how it was caught, and its effect.}
\label{tab:errors}
\begin{tabular}{p{3.4cm} p{3.1cm} p{3.4cm} p{4.2cm}}
\toprule
Error & How it was caught & Correction & Effect on results \\
\midrule
Calendar-grid volatility: returns computed across weekend rows carrying Friday's price, while annualising by $\sqrt{252}$ trading days & The design's own falsification tests failed (a fake 2019 ``ban'' showed a significant effect); investigating that failure exposed the grid & Weekday filter before any return/volatility computation; full rebuild and re-estimation & Overturned the falsification failures; every volatility number changed; an apparent GARCH result evaporated \\
\addlinespace
Mislabelled guar series (id 75 mixed guar gum into guar seed) & Correlation check of every spot series against its futures (0.04 vs.\ 0.99) & Switched to the correct ``Guar Seed'' series (id 413) & Removed a contaminated control \\
\addlinespace
Paddy volatility treated as market volatility despite MSP price-pinning (40.3\% exactly-flat returns) & Flat-run diagnostic across all series & Paddy excluded from primary analysis; wheat flagged & Placebo test became cleaner; headline sharpened \\
\addlinespace
Overstated precision: a pooled synthetic-DiD ``$z=-7.7$'' & Internal referee-style audit of the inference code & The statistic measured donor-swap stability, not sampling error; withdrawn. Valid placebo-based inference reported instead & Honest pooled inference is borderline ($z\approx-1.9$ once enough donors exist) \\
\addlinespace
Over-optimistic p-value ($p=0.008$) from large-sample clustering with only $\sim$10 clusters & Same audit & Few-cluster inference adopted: cluster-robust $t$ with $G-1$ degrees of freedom, plus wild-cluster bootstrap & Headline significance moved from ``strong'' to borderline (then to insignificant under better donors) \\
\addlinespace
Variance-ratio test statistic carried a spurious sample-size factor (companion efficiency analysis) & Monte-Carlo check during audit (statistic's null variance was wrong by $\sqrt{n}$) & Corrected to the textbook form; conclusions re-checked & Point estimates unchanged; a wrong ``test uninterpretable'' caveat was retracted \\
\addlinespace
Comparison-group composition: industrial/spice controls (castor, guar, cotton, jeera, turmeric) are not economically comparable to food staples & Anticipated in the pre-registered design critique; confirmed empirically when food-crop controls were added & Food-cereal donors added to the pool; oilseeds in progress & The headline effect halved (from $-20.7\%$ to $-9.8\%$) and lost aggregate significance - see Section~\ref{sec:progression} \\
\bottomrule
\end{tabular}
\end{table}

\paragraph{The calendar-grid error (the most important one).} My first corrected estimates
looked strong but \emph{failed their own pre-registered falsification tests}: a placebo
``ban'' placed in December 2019 - a date at which nothing happened - produced a
statistically significant effect ($-13.6\%$, $p=0.046$), and the pre-trend test rejected.
By the rules I had fixed in advance, that should have retired the design. Rather than
soften the rules, I investigated the failure and found the cause in the data, not the
design: the mandi series were seven-day calendar grids (28.5\% of one series' rows fell on
weekends), so ``daily returns'' were being computed across carried-forward weekend prices
while the annualisation assumed trading days. Fixing the grid (and excluding
price-support-censored paddy) flipped both falsification tests to passing. The lesson I
take from this: the falsification battery did exactly its job - it caught a measurement
bug I did not know I had.

\paragraph{The comparison-group error (the most consequential for the conclusion).} The
first credible estimates used the commodities whose futures kept trading as controls -
castor, guar, cotton, jeera, turmeric. Those are industrial, fibre, and spice crops; the
banned commodities are food staples. When I added economically comparable \emph{food}
controls (barley, maize, jowar, bajra), the estimated effect \emph{halved} and lost
statistical significance - because the food controls' volatility also declined over
2022-25, meaning part of the earlier ``effect'' was just the difference between food and
non-food markets. I treat this donor-sensitivity not as an embarrassment but as the
project's most informative single result: \emph{the estimate answers to the quality of its
counterfactual.}

% =====================================================================
\section{Methods, in plain language}
\label{sec:methods}

\paragraph{Realized volatility.} For each commodity-district, daily log price changes are
computed on trading days; their 30-day rolling standard deviation, annualised by
$\sqrt{252}$, is that market's ``realized volatility''; the panel aggregates it monthly.
The analysis works in logs of this measure, so effects read as percentage changes.

\paragraph{Difference-in-differences (DiD).} Compare the before-vs-after change in banned
commodities' volatility with the same change in comparison commodities:
\[
\ln(\mathrm{vol}_{it}) \;=\; \alpha_i + \lambda_t + \beta\,(\mathrm{banned}_i \times
\mathrm{post}_t) + \varepsilon_{it},
\]
with commodity-district fixed effects $\alpha_i$ (absorbing permanent differences across
markets) and month fixed effects $\lambda_t$ (absorbing anything common to all
commodities in a month - e.g.\ the 2022 Ukraine shock). $\beta$ converts to a percentage
effect as $e^{\beta}-1$. Because chana and mustard were suspended earlier than the rest,
per-commodity treatment dates are used in the commodity-level estimators, and the pooled
regression was checked against staggered dating (immaterial difference).

\paragraph{Honest uncertainty with few clusters.} Treatment is assigned at the
\emph{commodity} level, so there are effectively only 14 independent units (5 treated),
however many district rows exist. Conventional standard errors badly overstate precision
in this situation. Two corrections are used: a cluster-robust $t$-statistic referred to a
$t$ distribution with $G-1$ degrees of freedom, and a wild-cluster bootstrap (resampling
whole commodities with random sign flips under the null). Both are reported everywhere;
the naive p-value is shown only to demonstrate how over-optimistic it is.

\paragraph{Falsification battery (fixed before estimation).} (i)~\emph{Placebo date}: run
the same design on pre-ban data with a fake ban in December 2019; a trustworthy design
should find nothing. (ii)~\emph{Pre-trend test}: event-study ``leads'' should be jointly
zero before the ban. Failing either retires the design under the pre-registered rules.

\paragraph{Synthetic control (SCM).} For each banned commodity, construct a ``synthetic
twin'' - a weighted average of comparison commodities chosen to track the banned
commodity's pre-ban volatility path - and read the effect as the post-ban gap between
the commodity and its twin. Significance is assessed by in-space placebos (pretending
each comparison commodity was treated): with nine donors the attainable p-value floor is
$1/10 = 0.10$, a structural limitation reported alongside the estimates. A district-level
variant fits one synthetic control per treated district.

\paragraph{Synthetic difference-in-differences (SDID).} A hybrid that both reweights
comparison commodities (like SCM) and reweights pre-ban time periods, then takes the
double difference; uncertainty comes from placebo replications. With the enlarged donor
pool this estimator finally has a valid standard error.

\paragraph{Conditional volatility (GARCH) with break detection.} For the claim that
volatility \emph{dynamics} changed: rather than naively splitting at the ban date (which
biases persistence estimates when a variance break is mis-dated), structural variance
breaks are first located by an ICSS scan (with the heteroskedasticity-robust correction
appropriate for fat-tailed returns), and GARCH(1,1) is estimated within regimes; the
question is whether any detected break sits near the suspension date.

\paragraph{Companion designs.} (i)~\emph{Spot-market efficiency and integration}:
variance-ratio tests of return predictability and cross-mandi cointegration (do prices in
different markets move together, and how fast do gaps close), before vs.\ after, banned
vs.\ comparison. (ii)~\emph{Harvest troughs}: seasonal price indices testing whether
harvest-time price dips deepened once the forward price signal was gone.
(iii)~\emph{Market composition}: participant-wise turnover shares from regulator
bulletins, characterising who actually used these markets before the suspension.

% =====================================================================
\section{How the findings evolved}
\label{sec:progression}

The findings changed materially - twice - as the data and the comparison group
improved. I document the progression explicitly because each revision was forced by a
check working as intended, and because the direction of the revisions is itself
informative.

\begin{table}[htbp]
\centering\small
\caption{Progression of the headline volatility estimate.}
\label{tab:progression}
\begin{tabular}{p{3.1cm} p{3.3cm} p{3.4cm} p{4.0cm}}
\toprule
Stage & Headline DiD & Falsification & Verdict at that stage \\
\midrule
1. Initial hypothesis (from early exploratory work) & $+8$ to $+10\%$ (volatility \emph{rose}) & - & The claim this project set out to verify \\
\addlinespace
2. First estimates (calendar-grid data) & $-16.4\%$ ($p=0.029$) & \textbf{Failed}: placebo $-13.6\%$ ($p=0.046$); pre-trends rejected ($p=0.033$) & Design dead by its own rules; cause investigated \\
\addlinespace
3. After the trading-day fix and paddy exclusion & $-20.7\%$; few-cluster $p\approx0.026$ / bootstrap $0.045$ & Passed, but marginally: placebo $-11.7\%$ ($\approx$56\% of headline); a late pre-ban lead bin still hot & ``Volatility fell $\sim$20\%'' - borderline significant, industrial/spice controls \\
\addlinespace
4. Internal referee-style audit of code and inference & unchanged point estimate & - & Invalid pooled-SDID precision withdrawn; few-cluster inference adopted; leave-one-out added; donor pool flagged as the weak point \\
\addlinespace
5. Food-cereal donors added (current) & $-9.8\%$; $p\approx0.145$ / bootstrap $0.153$ (\textbf{not significant}) & \textbf{Clean}: placebo $-6.5\%$ ($p\approx0.34/0.43$); pre-trends $p\approx0.45$, no significant lead bin & Aggregate effect modest and insignificant; signal concentrated in chana; wheat's decline revealed as confounded \\
\bottomrule
\end{tabular}
\end{table}

\begin{keybox}
The two revisions moved the result in the same honest direction: from a strong claim on
weak foundations to a modest claim on solid ones. The initial ``$+8$-$10\%$ rise''
hypothesis is refuted at every stage. The ``$-20\%$ fall'' of stage~3 does not survive a
better comparison group. What survives everything: spot volatility did \emph{not} rise;
the measured decline is about $-10\%$ and statistically indistinguishable from zero at
conventional levels; and the one commodity with strong, consistent evidence of calming is
chana - the cleanest, least policy-contaminated banned commodity.
\end{keybox}

% =====================================================================
\section{Current results}
\label{sec:results}

\subsection{The volatility effect (headline analysis)}

\begin{table}[htbp]
\centering\small
\caption{Current estimates of the suspension's effect on spot realized volatility
(trading-day data; paddy excluded; 5 treated and 9 comparison commodities).}
\label{tab:results}
\begin{tabular}{p{4.3cm} p{2.9cm} p{6.4cm}}
\toprule
Estimator & Effect & Honest inference \\
\midrule
Two-way fixed-effects DiD & $-9.8\%$ & few-cluster $p=0.145$; wild bootstrap $p=0.153$ - \emph{not significant} \\
Synthetic DiD, pooled & $-19.1\%$ & placebo standard error valid with 9 donors; $z=-1.88$ (borderline) \\
Synthetic DiD, per commodity & chana $-38.7$; mustard $-15.2$; wheat $+0.4$; soybean $-26.1$; moong $-11.5$ (\%) & chana $z=-2.05$ (significant); all others insignificant; wheat $\approx$ zero \\
Synthetic control, per commodity & chana $-37.3$; mustard $-28.2$; wheat $-37.8$; soybean $-27.6$; moong $-8.7$ (\%) & chana at the placebo floor ($p=0.100$); others weak ($p=0.3$-$0.6$) \\
District-level synthetic control & negative for all five & chana/mustard/wheat at the $p=0.100$ floor; the floor itself ($1/10$) precludes 5\%-level claims \\
\bottomrule
\end{tabular}
\end{table}

Three readings of Table~\ref{tab:results} matter. First, the estimators agree on
\emph{sign} but their agreement is partly mechanical - they share one panel and one
treated set - so they are corroboration, not independent confirmations. Second, the
signal is \emph{concentrated in chana}: it is the only commodity individually significant
under any estimator, and the leave-one-out analysis below shows the aggregate thins to
$-7.4\%$ without it. Third, \emph{wheat's} apparent volatility decline collapses to zero
under the better donors - consistent with its decline having been driven by the export
ban, stock limits, and support-price operations that hit wheat simultaneously, not by the
futures suspension.

\begin{table}[htbp]
\centering\small
\caption{Falsification and robustness (current panel).}
\label{tab:falsification}
\begin{tabular}{p{5.2cm} p{4.2cm} p{4.2cm}}
\toprule
Check & Result & Reading \\
\midrule
Placebo ``ban'' in Dec.\ 2019 & $-6.5\%$, $p=0.344$ (bootstrap $0.433$) & design finds nothing where nothing happened \\
Joint pre-trend test & $p=0.445$; no individually significant lead bin & parallel-trends assumption not rejected \\
Leave-one-commodity-out & sign survives every drop ($-7.4\%$ to $-15.2\%$) & ex-wheat: $-15.2\%$ ($p\approx0.003$, significant); ex-chana: $-7.4\%$ ($p\approx0.29$) \\
District-liquidity / outlier filters & $-8.7\%$ to $-10.0\%$ & not an artifact of thin, noisy districts \\
\bottomrule
\end{tabular}
\end{table}

\subsection{Volatility dynamics (GARCH)}
The complementary claim that the ban shifted volatility \emph{dynamics} finds no support:
the break-detection scan locates no variance break near the suspension date in any
banned commodity's series (the only detected break in chana sits in September 2018), and
an early, apparently dramatic GARCH result - chana's volatility persistence falling
from 0.99 to 0.78 - turned out to be an artifact of the calendar-grid error and is
roughly flat on clean data. The volatility case therefore rests on the panel analysis
above, not on GARCH.

\subsection{Companion analyses: the cost side}
\label{sec:results-companions}

\begin{itemize}[leftmargin=1.6em,itemsep=2pt]
  \item \textbf{Who actually used these markets (market composition).} In the year before
        the suspension, the NCDEX agricultural segment's turnover was $\approx$42\%
        proprietary and $\approx$57\% client trading, with registered \emph{hedgers at
        only $\approx$1.4\%}. Direct exchange-based hedging was minimal even before the
        ban - so ``farmers lost their hedging tool'' is not where the cost mainly lies.
  \item \textbf{Spot-market efficiency and integration.} After the suspension, banned
        commodities' spot markets became \emph{less} informationally efficient
        (variance-ratio deviations grew relative to controls) and \emph{less} spatially
        integrated: the share of cointegrated market pairs fell by 0.22 more for banned
        than comparison commodities (one-sided $p\approx0.07$; only 4 banned vs.\ 5
        comparison commodities, so power is limited), and price gaps between markets
        closed more slowly. The comparison commodities meanwhile \emph{improved} -
        so this estimate is, if anything, conservative.
  \item \textbf{Harvest troughs.} Harvest-season price dips did \emph{not} deepen for
        banned commodities after the ban (the point estimate goes slightly the other way,
        $p=0.20$) - no evidence that losing the forward price signal worsened
        harvest-time price collapses, though the quantity (arrivals) data needed to test
        the mechanism directly are not available.
\end{itemize}

\begin{keybox}
Taken together: the suspension did not destabilise spot prices - and the evidence for
active \emph{calming} is modest and chana-specific - while the measurable cost shows up
in \emph{price discovery and market integration}, not in direct hedging (which barely
existed). ``Stability delivered, information lost'' is a fair one-line summary of where
the evidence currently points, with both halves held to honest standards of uncertainty.
\end{keybox}

% =====================================================================
\section{Interpretation, limitations, and next steps}
\label{sec:limits}

\paragraph{What can and cannot be claimed.} This is a quasi-experimental evaluation of an
\emph{endogenous} policy: the suspension was imposed \emph{because} prices had run up, so
some post-ban calming would be expected from mean reversion alone. The passing placebo
and pre-trend tests reduce that worry; they do not eliminate it. The defensible claims
are: (i)~the hypothesis that the ban \emph{raised} spot volatility is refuted;
(ii)~the aggregate volatility decline is about $-10\%$ and not statistically significant
under honest inference; (iii)~chana shows a strong, consistent, individually significant
decline; (iv)~wheat's decline is confounded by simultaneous interventions and should not
be attributed to the suspension.

\paragraph{Known limitations.}
\begin{itemize}[leftmargin=1.6em,itemsep=1pt]
  \item \emph{Few treated units.} Five treated commodities is an irreducible constraint;
        all inference here is honest about it, but no method fully escapes it.
  \item \emph{Incomplete donor set.} The oilseed comparison crops (groundnut, sesamum,
        sunflower) and ragi are still downloading; mustard and soybean results are
        provisional until they are in.
  \item \emph{Bundled treatment.} Export bans, stock limits, and support-price operations
        hit several banned commodities in the same window; commodity-specific attribution
        is weakest for wheat (and is why chana - least contaminated - carries the
        evidentiary weight).
  \item \emph{Missing banned-commodity futures.} The pre-ban NCDEX futures for six banned
        commodities are unavailable pending an exchange data request; price-discovery and
        basis analyses that need them are specified but not executed.
  \item \emph{Planned upgrades not yet run.} A formal pre-trend sensitivity analysis
        (Rambachan-Roth style), a modern staggered-adoption estimator
        (Callaway-Sant'Anna), and a bootstrap confidence interval for the headline are
        specified as next steps.
\end{itemize}

% =====================================================================
\section{Reproducibility and folder navigation}
\label{sec:repro}

Everything in this report can be regenerated from the repository. The layout, the key
result files, and the exact commands are documented in \texttt{FOLDER\_GUIDE.md} at the
repository root; in brief:

\begin{itemize}[leftmargin=1.6em,itemsep=1pt]
  \item \texttt{02\_data/} - raw (immutable), constructed, and clean data tiers;
        \texttt{00\_admin/data\_log.csv} records the provenance of every series.
  \item \texttt{00\_code/} - all acquisition, cleaning, construction, and estimation
        scripts; each estimation script writes its outputs to the matching hypothesis
        folder under \texttt{04\_empirics/}.
  \item \texttt{04\_empirics/H1\_volatility/c1\_findings.md} - the canonical results
        memo (kept in lock-step with the result files); \texttt{output/} folders hold
        every regression output, placebo, bootstrap, and figure cited here.
  \item \texttt{05\_paper/drafts/paper.tex} - the academic paper draft;
        \texttt{make\_paper\_tables.py} regenerates its tables directly from the result
        files, so paper numbers cannot drift from outputs.
  \item \texttt{03\_policy\_timeline/policy\_interventions.csv} - the confounder
        ledger, one primary-sourced row per policy event.
\end{itemize}

The provenance of every acquired data series - source, access date, coverage, units,
and known issues - is recorded one row per series in \texttt{00\_admin/data\_log.csv};
the methodological history (what was decided when, which errors were found, and how the
estimates evolved) is documented in Sections~\ref{sec:errors} and~\ref{sec:progression}
of this report.

\end{document}
